\documentclass[14pt,a4paper]{extarticle}
\usepackage[T2A]{fontenc}
\usepackage[utf8]{inputenc}
\usepackage[english,russian]{babel}
\usepackage{lmodern}
\usepackage{geometry}
\usepackage{array}
\usepackage{tabularx}
\usepackage{longtable}
\usepackage{ltablex}
\usepackage{booktabs}
\usepackage{hyperref}
\usepackage{indentfirst}
\usepackage{microtype}
\usepackage{graphicx}
\usepackage{setspace}
\usepackage{listings}
\usepackage{graphicx}
\usepackage{caption}
\usepackage{float}
\usepackage{tocloft}
\usepackage{titlesec}
\usepackage{newtxtext,newtxmath}

\titleformat{\section}        % форматируем \section
  {\normalfont\bfseries}      % стиль: обычный + жирный
  {\thesection.}              % номер секции с точкой
  {2mm}                       % расстояние между номером и текстом
  {}                           % перед заголовком можно добавить что-то
\titlespacing*{\section}{1.5cm}{*3}{*1}  % {отступ слева}{отступ сверху}{отступ снизу}

\captionsetup{justification=centering}
\renewcommand{\ttdefault}{cmtt}
\keepXColumns
\setlength{\parindent}{3em}
\newgeometry{left=2.5cm,right=2.5cm,top=3cm,bottom=2.5cm}

\renewcommand{\contentsname}{\hspace*{4em}{\smallСодержание}}
\setlength{\cftbeforesecskip}{0.5em}
\renewcommand{\cftsecleader}{\cftdotfill{\cftdotsep}}
\renewcommand{\cftsecaftersnum}{.\hspace{2mm}}

\title{Отчёт по "Организация работы ИТ-службы"}
\author{Выполнил студент группы М-РИТ-25-1, Станиславчук Сергей}

\begin{document}

\setstretch{1.5}
\noindent
\begin{minipage}[t]{0.01\textwidth}
    \includegraphics{~/programmer/++Programmer/5_1/LSTU.jpg}
\end{minipage}%
\begin{minipage}[c]{1.1\linewidth}
    \centering
    {\small
    МИНИСТЕРСТВО НАУКИ И ВЫСШЕГО ОБРАЗОВАНИЯ РФ \\[1em]
    ФЕДЕРАЛЬНОЕ ГОСУДАРСТВЕННОЕ БЮДЖЕТНОЕ \\
    ОБРАЗОВАТЕЛЬНОЕ УЧРЕЖДЕНИЕ ВЫСШЕГО ОБРАЗОВАНИЯ \\
    «ЛИПЕЦКИЙ ГОСУДАРСТВЕННЫЙ ТЕХНИЧЕСКИЙ УНИВЕРСИТЕТ»
    }
\end{minipage}%

\vspace{3em}

\noindent
\hspace{1em}%
\begin{minipage}[t]{0.8\textwidth}
\textbf{Институт} \hspace{2em} компьютерных наук \\
\textbf{Кафедра} \hspace{2em} автоматизированных систем управления
\end{minipage}


\vspace{6em}

\begin{center}
\textbf{Практическая работа №3} \\
по дисциплине \\ <<Организация работы ИТ-службы>>
\end{center}

\vspace{6em}

\noindent
\begin{tabular*}{\textwidth}{@{\extracolsep{\fill}} l c r }
    Студент М-РИТ-25-1 & \makebox[4cm]{\hrulefill} & Станиславчук С.М. \\[2em]
Руководитель & \\
К.т.н., доцент & \makebox[4cm]{\hrulefill} & Ведищев В.В. \\
\end{tabular*}

\vfill
\begin{center}
    Липецк, 2026
\end{center}

\thispagestyle{empty}

\newpage
\tableofcontents
\newpage

\section{Новые или изменяемые услуги}

Проектирование новых или изменяемых услуг направлено на обеспечение соответствия IT-сервиса потребностям заказчиков и требованиям бизнеса. Изменения могут быть связаны с появлением новых услуг, расширением функциональности существующих предложений, модернизацией инфраструктуры или изменением требований к качеству обслуживания.

В случае VDS-хостинга примером изменения услуги может быть расширение линейки конфигураций виртуальных серверов, увеличение доступного дискового пространства или внедрение дополнительных возможностей хранения данных. Необходимость изменений может возникнуть вследствие увеличения спроса, появления новых технологий или изменения требований заказчиков к производительности и доступности.

Перед внедрением изменения необходимо определить его цели, оценить затраты, технические риски и влияние на действующих клиентов. Например, модернизация платформы виртуализации может повысить производительность, но одновременно потребовать переноса виртуальных машин и временных ограничений обслуживания.

Таким образом, проектирование новых и изменяемых услуг должно обеспечивать согласование потребностей заказчиков с техническими и экономическими возможностями поставщика.

\section{Системы управления услугами и инструментарий, портфель и каталог услуг}

Для эффективного предоставления VDS-хостинга необходимо использовать инструменты управления услугами, позволяющие контролировать состояние инфраструктуры, обрабатывать обращения клиентов и управлять изменениями.

К таким инструментам относятся системы мониторинга, средства автоматизированного создания виртуальных серверов, системы регистрации обращений, инструменты управления конфигурациями и средства резервного копирования. Они позволяют сократить время выполнения типовых операций, своевременно обнаруживать неисправности и поддерживать установленный уровень качества обслуживания.

Портфель услуг отражает совокупность услуг поставщика на разных этапах жизненного цикла, включая разрабатываемые, действующие и выводимые из эксплуатации услуги. Для VDS-услуги портфель может включать основные предложения по аренде виртуальных серверов и дополнительные услуги, связанные с хранением данных и управлением ресурсами.

Каталог услуг представляет собой структурированное описание услуг, доступных заказчикам. Для VDS-хостинга он должен содержать описание доступных конфигураций, объём процессорных ресурсов и оперативной памяти, размер дискового пространства, параметры сетевого подключения, стоимость и условия предоставления услуги.

Портфель используется для стратегического управления составом услуг, тогда как каталог помогает заказчикам выбирать и заказывать конкретные предложения. Совместное использование этих инструментов упрощает управление услугами и обеспечивает прозрачность условий их предоставления.

\section{ИТ-инфраструктура, её архитектура и системы}

ИТ-инфраструктура является технической основой VDS-хостинга. Она включает физическое оборудование, программное обеспечение, системы хранения данных, сетевую инфраструктуру и средства управления виртуализацией.

На физическом уровне располагаются серверы, обеспечивающие вычислительные мощности, системы хранения данных, предназначенные для постоянного хранения виртуальных дисков и данных клиентов, а также сетевое оборудование, обеспечивающее взаимодействие между компонентами инфраструктуры и доступ клиентов к виртуальным серверам.

На программном уровне используется платформа виртуализации, позволяющая создавать и запускать виртуальные машины на физических серверах. Каждая виртуальная машина предоставляет заказчику изолированную вычислительную среду, в которой могут устанавливаться операционная система и необходимые приложения.

Для управления инфраструктурой необходимы системы мониторинга, автоматизации, резервного копирования и информационной безопасности. Мониторинг позволяет контролировать загрузку процессоров, использование оперативной памяти, состояние хранилищ и доступность сетевых соединений. Резервное копирование и процедуры восстановления помогают снизить риск потери данных при отказах.

Архитектура инфраструктуры должна учитывать масштабируемость, производительность, безопасность и устойчивость к отказам. Например, при увеличении количества клиентов необходимо предусмотреть возможность добавления физических серверов и расширения систем хранения данных без существенного нарушения предоставления действующих услуг.

\section{Требования к процессам управления}

Проектирование VDS-хостинга должно учитывать не только технические характеристики инфраструктуры, но и требования к процессам управления, обеспечивающим стабильное предоставление услуги.

Процесс управления инцидентами должен обеспечивать регистрацию, классификацию и устранение нарушений нормальной работы сервиса. Например, при недоступности виртуального сервера необходимо определить причину сбоя, восстановить работоспособность и уведомить затронутого клиента.

Управление изменениями должно обеспечивать контролируемое внесение изменений в инфраструктуру. Перед обновлением платформы виртуализации или заменой оборудования необходимо оценить риски, определить возможное влияние на доступность услуги и спланировать выполнение работ.

Управление запросами на обслуживание должно регламентировать выполнение типовых операций по обращениям клиентов. К ним относятся изменение конфигурации виртуального сервера, увеличение дискового пространства и другие предусмотренные поставщиком операции.

Управление доступностью и мощностями должно обеспечивать соответствие возможностей инфраструктуры установленным требованиям. Для этого необходимо контролировать загрузку физических серверов, прогнозировать потребность в ресурсах и планировать расширение инфраструктуры.

Отдельное значение имеет управление информационной безопасностью. Процессы должны предусматривать контроль доступа, защиту инфраструктуры от несанкционированного вмешательства, обработку сообщений о нарушениях и соблюдение правил использования сервиса.

Для каждого процесса необходимо определить ответственных лиц, порядок выполнения операций, используемые инструменты и критерии оценки результатов. Регламентация процессов позволяет уменьшить вероятность ошибок и обеспечить предсказуемое качество обслуживания.

\section{Методы измерения и метрики}

Методы измерения и метрики используются для оценки качества VDS-хостинга, контроля выполнения требований и выявления направлений улучшения. Измерения должны охватывать как техническое состояние инфраструктуры, так и качество процессов обслуживания клиентов.

Одним из основных показателей является доступность сервиса. Она определяется отношением времени, в течение которого услуга была доступна, к общему согласованному времени обслуживания. Например, при целевой доступности 99,9\% максимально допустимое время недоступности составляет приблизительно 43,8 минуты за 30-дневный период, если применяется соответствующая методика расчёта и не учитываются исключения, предусмотренные SLA.

Производительность оценивается по загрузке процессора, использованию оперативной памяти, задержкам дисковых операций, пропускной способности сети и времени отклика размещённых приложений. Эти показатели позволяют выявлять недостаток ресурсов и определять необходимость изменения конфигурации или расширения инфраструктуры.

Для оценки надежности могут использоваться среднее время между отказами (MTBF) и среднее время восстановления (MTTR). Первый показатель характеризует промежуток времени между отказами, а второй — среднюю продолжительность восстановления после возникновения неисправности.

Качество технической поддержки может оцениваться по среднему времени реакции на обращение, среднему времени разрешения инцидента, количеству повторных инцидентов и доле обращений, обработанных в установленные сроки.

Экономическая эффективность оценивается с помощью затрат на предоставление услуги, загрузки доступных ресурсов и доходов от обслуживания клиентов. Анализ этих показателей помогает определить, насколько рационально используются ресурсы и требуется ли корректировка тарифов или инфраструктуры.

Сбор и анализ метрик позволяют сопоставлять фактическое качество сервиса с установленными целями. На основании результатов измерений принимаются решения о модернизации инфраструктуры, изменении процессов управления и дальнейшем развитии VDS-хостинга.

\end{document}
